% !TEX program = xelatex
% 本卷为范围内改编题，无对应官方解析原页；仅作原扫描逐字核验与数学批注。
\documentclass[11pt,a4paper,fontset=fandol]{ctexart}
\usepackage[a4paper,left=20mm,right=20mm,top=20mm,bottom=18mm,headheight=15pt]{geometry}
\usepackage{amsmath,amssymb,fancyhdr,lastpage,xcolor,cancel}
\definecolor{accent}{HTML}{2C4A7C}\definecolor{problemred}{HTML}{B42318}\definecolor{fixgreen}{HTML}{137333}
\setlength{\parindent}{0pt}\setlength{\parskip}{0.45em}\linespread{1.12}
\pagestyle{fancy}\fancyhf{}\fancyhead[L]{\small 数学二}\fancyhead[R]{\small 第3讲收尾检测}\fancyfoot[C]{\small 第 \thepage\ 页，共 \pageref{LastPage} 页}\renewcommand{\headrulewidth}{0.6pt}
\newcommand{\pt}[1]{\par\vspace{.3em}{\large\bfseries\color{accent}#1}\par\vspace{.1em}}
\newcommand{\problem}[1]{\textcolor{problemred}{\textbf{问题：}#1}}
\newcommand{\fix}[1]{\textcolor{fixgreen}{\textbf{修正：}#1}}
\begin{document}
\section*{第1题}\pt{原题}设函数 $f$ 在 $0$ 处可导，且 $f(0)=0$。求 $\displaystyle\lim_{x\to0}\frac{f(\sin x)}x$。
\pt{我的解答}
$f$在$0$处可导，$f(0)=0$，求$\displaystyle\lim_{x\to0}\frac{f(\sin x)}x$\\
$\displaystyle\lim_{x\to0}\frac{f(\sin x)}x=\lim_{x\to0}\frac{f'(\sin x)\cdot\cos x}{\cancel{x}}$。\quad\problem{这里的分母 $x$ 已被你划掉；按有效书写，应读作洛必达后分母求导为 $1$，不能再把它转录成保留 $x$ 的式子。}\\
$=f'(\cos x)$
\pt{正确解答}
\fix{按当前题设的导数定义：}$\displaystyle\frac{f(\sin x)}x=\frac{f(\sin x)-f(0)}{\sin x-0}\cdot\frac{\sin x}{x}\to f'(0)\cdot1=f'(0)$。

\section*{第2题}\pt{原题}设 $\displaystyle f(x)=\begin{cases}ax^2+bx+1,&x\le1,\\\ln x+c(x-1),&x>1,\end{cases}$ 若 $f$ 在 $x=1$ 处可导，求 $a,b,c$ 之间应满足的全部条件。
\pt{我的解答}
$f(x)=\begin{cases}ax^2+bx+1,&x\le1\\\ln x+c(x-1),&x>1\end{cases}$，$f$在$x=1$处可导。\\
$f(1)=a+b+1$\\
$f_-(1)=\lim_{x\to1^-}f(x)=a+b+1$\\
$f_+(1)=\lim_{x\to1^+}f(x)=0$\\
$\Rightarrow a+b=-1$\\
$f'_-(1)=\lim_{x\to1^-}\frac{f(1)-f(x)}{1-x}=\lim_{x\to1^-}\frac{ax^2+bx+1}{x-1}=\lim_{x\to1^-}(2ax+b)=2a+b$。\quad\problem{纸面这里的分式中，分子处有划改；有效的左导数定义应保留 $f(1)-f(x)$，再把分子化为 $f(1)-(ax^2+bx+1)$，不能直接写成 $ax^2+bx+1$。}\\
$f'_+(1)=\lim_{x\to1^+}\frac{f(x)-f(1)}{x-1}=\lim_{x\to1^+}\frac{\ln x+c(x-1)}{x-1}=\lim_{x\to1^+}(\frac1x+c)=c+1$\\
$\Rightarrow2a+b=c+1$
\pt{正确解答}
\fix{连续性：}$a+b+1=0$，即 $a+b=-1$。\quad
\fix{左右导数：}$2a+b=c+1$。因此 $\boxed{a+b=-1,\ 2a+b=c+1}$，等价于 $b=-a-1,\ c=a-2$。

\section*{第3题}\pt{原题}设 $\displaystyle f(x)=\begin{cases}\frac{\sin(ax)}x,&x\ne0,\\b,&x=0,\end{cases}$ 其中 $a,b$ 为常数。求使 $f$ 在 $x=0$ 处可导的 $a,b$。
\pt{我的解答}
$f(x)=\begin{cases}\dfrac{\sin(ax)}x,&x\ne0\\b,&x=0\end{cases}$，求$f$在$x=0$处可导。\\
$f(0)=b$，$\lim_{x\to0}f(x)=\lim_{x\to0}\dfrac{\sin(ax)}x=a$\\
$\Rightarrow a=b$\\
$\therefore f'(0)=\lim_{x\to0}\dfrac{f(x)-f(0)}x=\lim_{x\to0}\dfrac{\frac{\sin(ax)}x-b}{x}=\lim_{x\to0}\dfrac{\sin(ax)-bx}{x^2}$\\
$=\lim_{x\to0}\dfrac{a\cos(ax)-b}{2x}$。\quad\problem{到这里使用了一次洛必达；代入 $a=b$ 后仍是 $0/0$，不能仅由分子趋于 $0$ 得出极限值。}\\
$\Rightarrow x\to0$时，$\cos(ax)\cdot a-b\to0$\\
$\Rightarrow a-b\to0，a=b$
\pt{正确解答}
\fix{保留 $a=b$；补足极限：}当 $b=a$ 时，
\[
f'(0)=\lim_{x\to0}\frac{\sin(ax)-ax}{x^2}
=\lim_{x\to0}\frac{a\cos(ax)-a}{2x}
=\frac a2\lim_{x\to0}\frac{\cos(ax)-1}{x}.
\]
又因为
\[
\frac{1-\cos(ax)}x
=2\frac{\sin^2(ax/2)}x
=\frac{a^2x}{2}\left(\frac{\sin(ax/2)}{ax/2}\right)^2\to0,
\]
所以 $f'(0)=0$。因此全部条件为 $\boxed{a=b}$，其中 $a$ 任意。

\section*{第4题}\pt{原题}曲线 $y=f(x)$ 在点 $(1,2)$ 处有切线。若 $\displaystyle\lim_{h\to0}\frac{f(1+2h)-f(1)}h=6$，求该切线方程。
\pt{我的解答}
$\lim_{h\to0}\dfrac{f(1+2h)-f(1)}h=6$，求$y=f(x)$在$(1,2)$切线方程\\
$f(1)=2$\\
$\lim_{h\to0}\dfrac{f(1+2h)-2}{h}=6$\\
$\Rightarrow\lim_{h\to0}\dfrac{f(1+2h)-2}{2h}=3$\\
$\Rightarrow f'(1)=\lim_{2h\to0}\dfrac{f(1+2h)-f(1)}{2h-0}=3$\\
$\Rightarrow y-2=3(x-1)$
\end{document}
